En una regió de l'espai hi ha un camp magnètic constant dirigit cap a l'interior del paper. En aquesta regió entren dos electrons amb la mateixa rapidesa i la mateixa direcció, però movent-se en sentits contraris, tal com indica la figura.

\figura[0.4]{fig1.png}

\begin{apartat}{1}
Dibuixeu la força magnètica que actua sobre cada electró quan entra en la regió on hi ha el camp magnètic. Justifiqueu i dibuixeu les trajectòries dels dos electrons i indiqueu el sentit de gir.
\end{apartat}

\begin{apartat}{1}
Eliminem aquest camp magnètic i el substituïm per un altre camp magnètic, de manera que els electrons no es desvien quan entren en aquesta regió. Dibuixeu com hauria de ser aquest nou camp magnètic. Justifiqueu la resposta.
\end{apartat}

\nota{No és vàlida la resposta $\vec{B} = 0$.}
